\centerline{\bf \Huge Разнобой}
\centerline{\bf \Huge }

\begin{flushright}
        \scriptsize{}
\end{flushright}

\problem{1}
В лыжной гонке участвовало 50 спортсменов. Они стартовали друг за другом с одинаковыми интервалами, каждый со своей постоянной скоростью. Известно, что каждый лыжник лидировал в некоторый момент гонки. На каком месте финишировал спортсмен стартовавший десятым?

\problem{2}
2. Несколько человек построились в колонну, в том числе Андрей, Боря и Вова. Известно, что:

- людей, стоящих впереди Андрея, в 2 раза меньше, чем людей стоящих позади него;

- людей, стоящих впереди Бори, в 3 раза меньше, чем людей стоящих позади него;

- людей, стоящих впереди Вовы, в 4 раза меньше, чем людей стоящих позади него.

Какое наименьшее количество людей может быть в колонне?

\problem{3}
У Ани есть карточки с числами от 1 до 40. Она расставляет их в некотором порядке по кругу. После этого мама считает 40 разностей двух соседних чисел (из большего числа вычитает меньшее) и находит среди них наименьшую. После этого мама даёт Ане столько конфет, чему равна данная разность. Какое наибольшее количество конфет может себе гарантировать Аня?

\problem{4}
4. В четырёхугольнике $A B C D\ A D \| B C$ и $\angle A B C=2 \angle A D C$. На луче $B A$ за точкой $A$ отметили точку $P$ такую, что $A P=B C$.
Докажите, что $A D=A B+B C$ и треугольник CPD равнобедренный.

\problem{5}
В ряд положили 75 монет, какие-то орлом вверх, остальные - решкой. Известно, что любые две монеты между которыми ровно 5 других монет лежат по-разному (одна решкой вверх, а другая - орлом). Какое наибольшее количество монет может лежать орлом вверх?